\documentclass[12pt]{article}
\usepackage[utf8]{inputenc}
\usepackage[T1]{fontenc}
\usepackage{hyperref}
\title{Sahu Language Resources for K-12 Curriculum Development}
\author{Compiled by Claude Sonnet 5.5 (claude-sonnet-5-5) for Andr\'as Kornai}
\date{}
\begin{document}
\maketitle

\section{Introduction}
Sahu (ISO 639-3 saj; Glottocode sahu1245) is a North Halmahera language of Indonesia. Its dialects include Pa'disua, Tala'i, Waioli and Gamkonora. The searches found one major scholarly description and some recent and archival work, but no school materials.

\section{Primary Educational Resources}
No primers, textbooks or school programmes for Sahu were found in the search results.

\section{Linguistic and Reference Materials}
\begin{itemize}
\item Visser, L. E. and Voorhoeve, C. L. (1987), \emph{Sahu-Indonesian-English Dictionary and Sahu Grammar Sketch}, Foris, Dordrecht (KITLV Verhandelingen 126, 276 pages), Glottolog reference entry: \url{https://glottolog.org/resource/reference/id/87812}
\item Wikipedia article on the Sahu language: \url{https://en.wikipedia.org/wiki/Sahu_language}
\item ANU Open Research repository item surfaced in the search (content not inspected, so its relevance is unconfirmed): \url{https://openresearch-repository.anu.edu.au/server/api/core/bitstreams/5e57480c-3473-448c-9cca-5f1c6f98b7b7/content}
\end{itemize}

\section{Religious and Community Resources}
No scripture or community resources were found in these searches.

\section{Digital and Audio Resources}
\begin{itemize}
\item Paper (2024) on an Android bilingual dictionary for Sahu Tala'i and Indonesian with audio pronunciation and example sentences, built with the Aho-Corasick algorithm: \url{https://journal.mmtc.ac.id/index.php/pekommas/article/view/5281}
\item Max Planck Institute for Psycholinguistics, The Halmahera Collection (1980s recordings of Sahu speech and songs by Leontine Visser): \url{https://www.mpi.nl/de/node/53801}
\item PARADISEC series INHALM (Halmahera, Moluccas) field notes: \url{https://paradisec.org.au/fieldnotes/INHALM.htm}
\end{itemize}

\section{Curriculum Development Context}
Sahu is spoken by a small community and is not known to be taught in schools. The Visser and Voorhoeve dictionary and the archival recordings are the strongest foundation for any materials, and the Android dictionary shows recent local interest in the language. Community consultation and the status of Indonesian-medium schooling would need to be established before a curriculum is designed.

\end{document}